\documentclass[11pt,a4paper]{article}
\usepackage[T1]{fontenc}
\usepackage[utf8]{inputenc}
\usepackage{lmodern}
\usepackage{amsmath,amssymb,amsthm}
\usepackage{booktabs}
\usepackage{xcolor}
\usepackage[margin=27mm]{geometry}
\usepackage[colorlinks=true,urlcolor=blue,citecolor=blue,linkcolor=blue]{hyperref}
\newtheorem{theorem}{Theorem}
\newtheorem{proposition}[theorem]{Proposition}
\newtheorem{lemma}[theorem]{Lemma}
\newtheorem{corollary}[theorem]{Corollary}
\theoremstyle{definition}
\newtheorem{definition}[theorem]{Definition}
\newtheorem{remark}[theorem]{Remark}
\theoremstyle{plain}
\newcommand{\F}{\mathbb{F}_2}
\newcommand{\supp}{\operatorname{supp}}
\newcommand{\one}{\mathbf{1}}
\newcommand{\ket}[1]{\lvert#1\rangle}
\newcommand{\VV}{\mathrm{VV}}
\newcommand{\CC}{\mathrm{CC}}
\setlength{\emergencystretch}{2em}
\setlength{\parindent}{0pt}
\setlength{\parskip}{0.3em plus 0.1em minus 0.1em}
\title{A parity obstruction to depth-optimal syndrome extraction\\
for lifted-product codes}
\author{Fable \& Astra\\[0.35em]
\normalsize Research led by Jean-Philippe Bruneton\\[-0.1em]
\small Universit\'e Paris Cit\'e, France}
\date{19 November 2026}

\begin{document}
\maketitle
\begin{abstract}
Quantum error-correcting codes are run by repeatedly measuring their check
operators, each through an ancilla qubit that interacts with the data qubits
of its check by a sequence of controlled-NOT gates. The number of layers of
these gates sets the duration of an error-correction cycle, and it can never
be smaller than the largest number of gates that a single qubit must take
part in. Nguyen, Rimbach-Russ and Bosco recently constructed schedules for
lifted-product codes that reach this bound, or exceed it by one layer when
both factors of the code have odd degree, and left open whether the extra
layer is needed. We show that it is, for a broad class of codes: for the
lifted product of two $1\times1$ protographs whose entries have odd weights
$a$ and $b$, over a group of odd order, no valid schedule has depth
$\Delta=a+b$, so their depth $\Delta+1$ is optimal. The proof is a parity
count: summing the ordering conditions that make the circuit measure the
correct operators over all pairs of checks gives the odd number
$|G|\,ab$, while each condition requires an even contribution. The class
contains the $[[90,8,10]]$ bivariate bicycle code, for which the depth-7
circuit of Bravyi et al.\ is therefore optimal. Without the odd-order
hypothesis the obstruction can fail: an exhaustive search
finds a depth-$\Delta$ schedule for a $[[12,4]]$ code over the cyclic group of
order six.
\end{abstract}

\section{Introduction}
A stabilizer code of the Calderbank--Shor--Steane type is operated by
measuring its $X$-type and $Z$-type check operators again and again. In the
standard circuit each check owns an ancilla qubit, prepared in $\ket{+}$ for an
$X$ check and in $\ket{0}$ for a $Z$ check, which interacts with every data
qubit of its support through a controlled-NOT gate and is then measured. The
gates are arranged in layers in which each qubit takes part in at most one
gate. The depth, the number of layers, sets the duration of a cycle and the
amount of idle noise, so a short schedule matters for every hardware
platform. A qubit involved in $\Delta$ gates needs $\Delta$ distinct layers,
so the depth is at least the maximal degree $\Delta$ of the Tanner graph.

Not every assignment of layers measures the intended operators. When an $X$
check and a $Z$ check overlap, the order in which their two ancillas visit the
shared data qubits matters: the circuit is correct for this pair exactly when
the $X$ ancilla acts first on an even number of shared qubits
\cite{NRB2026}. For the surface code this is the familiar reason why the
order of the four gates around a plaquette cannot be chosen freely. These
\emph{properness} conditions are global and parity-like, and they are the
reason why the degree bound $\Delta$ is not always attained.

Nguyen, Rimbach-Russ and Bosco \cite{NRB2026} attacked this problem by
quotienting the Tanner graph of a structured code by its symmetries, solving
a much smaller scheduling problem, and lifting the solution back. For
lifted-product codes built from protographs $A$ and $B$ with maximal degrees
$\Delta_A$ and $\Delta_B$, their three-band schedule has depth
$\Delta_A+2\lceil\Delta_B/2\rceil$, which equals the lower bound
$\Delta=\Delta_A+\Delta_B$ when one of the two degrees is even and exceeds it
by one when both are odd. For the odd case they report that extensive
numerics on the Gross code suggest that $\Delta$ cannot be reached, but they
do not prove it.

The simplest instance shows where the obstruction comes from. Take $a=b=1$,
so that each check has weight two and each data qubit takes part in one $X$
and one $Z$ gate. The Tanner graph is then a disjoint union of four-cycles
$X$--data--$Z$--data, and a schedule of depth two must alternate its two
layers around each cycle. The $X$ ancilla then precedes the $Z$ ancilla on
exactly one of the two shared qubits, an odd number, and the circuit is
improper. Our results show that this one-cycle parity defect persists in
general, as a global count.

The main result is the following. Let $G$ be a finite group of odd order and
let $a,b\in\F[G]$ have odd weights, also denoted $a$ and $b$. Then no proper
schedule of the lifted product of the $1\times1$ protographs $(a)$ and $(b)$
has depth $\Delta=a+b$ (Theorem~\ref{thm:main}). The three-band schedule of
\cite{NRB2026}, of depth $\Delta+1$, is therefore optimal. The argument sums
the properness conditions over all pairs of checks: by double counting, the
total equals $|G|\,ab$ for every schedule of depth $\Delta$, which is odd under
our hypotheses, whereas properness makes each summand even. A second argument
(Proposition~\ref{prop:uniform}) shows that for every group, of any order,
no schedule that is uniform under the group action, which is the class in
which \cite{NRB2026} search, reaches depth $\Delta$. Both arguments need the
odd weights; the first also needs the odd order, and an explicit
depth-$\Delta$ schedule over the cyclic group of order six shows that this
hypothesis cannot simply be removed (Section~\ref{sec:even}).

\section{Related work}\label{sec:related}
Hypergraph-product codes were introduced by Tillich and Z\'emor
\cite{TillichZemor2014} and lifted-product codes by Panteleev and Kalachev
\cite{PanteleevKalachev2021}, who replaced the field of coefficients by a
group algebra. The lifted product of two $1\times1$ protographs is a
two-block group-algebra code in the sense of Lin and Pryadko
\cite{LinPryadko2023}; for an abelian group of the form
$\mathbb{Z}_\ell\times\mathbb{Z}_m$ these are the bivariate bicycle codes of
Bravyi et al.\ \cite{Bravyi2024}, whose syndrome cycle uses seven layers of
controlled-NOT gates for checks of weight six. Tremblay, Delfosse and
Beverland gave a general coloration circuit that measures the $X$ checks and
then the $Z$ checks, each in a number of gate layers equal to the maximal
degree of the corresponding Tanner graph \cite{TDB2022}; it does not
interleave the two types and so does not face the ordering conditions
studied here.

The present note answers the question left open by Nguyen, Rimbach-Russ and
Bosco \cite{NRB2026} for $1\times1$ protographs over groups of odd order. We
use their circuit model and their properness condition unchanged, and their
three-band construction supplies the matching upper bound. The Gross code
$[[144,12,12]]$ of \cite{Bravyi2024}, on which their numerics were run, is
built over a group of order $72$ and is not covered by our theorem, but the
$[[90,8,10]]$ code of the same family, over a group of order $45$, is.

\section{Model}\label{sec:model}
A CSS code is given by its Tanner graph $T$: a bipartite graph between the
check ancillas, each of type $X$ or $Z$, and the data qubits, with an edge for
each controlled-NOT gate. A \emph{schedule} of depth $L$ is a map
$\tau:E(T)\to\{1,\dots,L\}$ such that edges sharing a vertex receive distinct
layers. For an $X$ ancilla $x$, a $Z$ ancilla $z$ and a data qubit $q$ adjacent
to both, write $\tau(x,q)$ and $\tau(z,q)$ for the layers of the two gates. The
schedule is \emph{proper} if, for every pair $(x,z)$,
\begin{equation}\label{eq:proper}
\sum_{q\in N(x)\cap N(z)}\one\bigl[\tau(x,q)<\tau(z,q)\bigr]\equiv0\pmod 2,
\end{equation}
where $N(v)$ is the set of neighbours of $v$. This is condition~(2) of
\cite{NRB2026}. The overlaps $\lvert N(x)\cap N(z)\rvert$ are even because the
checks commute, so \eqref{eq:proper} can equivalently be stated with the
reverse order. We count only gate layers; preparation and measurement of the
ancillas are not included, and each check is measured by a single ancilla.

Let $G$ be a finite group of order $g$, and let $a,b\in\F[G]$ with supports
$\supp a$ and $\supp b$ of sizes $\lvert a\rvert$ and $\lvert b\rvert$. The lifted
product of the $1\times1$ protographs $(a)$ and $(b)$ has parity-check
matrices $H_X=[\,L(a)\;\;R(b)\,]$ and $H_Z=[\,R(b^*)\;\;L(a^*)\,]$, where $L$
and $R$ are the left and right regular representations and $^*$ is the
antipode $h\mapsto h^{-1}$ \cite{NRB2026}. Explicitly, there are $X$ ancillas
$X_u$, $Z$ ancillas $Z_u$ and data qubits $\VV_u$, $\CC_u$, indexed by $u\in G$,
with edges
\begin{equation}\label{eq:edges}
X_u - \VV_{hu},\quad X_u - \CC_{uk},\quad
Z_u - \VV_{uk^{-1}},\quad Z_u - \CC_{h^{-1}u},
\qquad h\in\supp a,\ k\in\supp b .
\end{equation}
The ancillas $X_u$ and $Z_{u'}$ share the qubit $\VV_{hu}$ when $hu=u'k^{-1}$
and the qubit $\CC_{uk}$ when $uk=h^{-1}u'$, that is, in both cases, for the
pairs $(h,k)\in\supp a\times\supp b$ with $u'=huk$. Each such pair gives one
shared $\VV$ qubit and one shared $\CC$ qubit, so the overlap is even and the
checks commute. Every ancilla has
degree $\Delta=\lvert a\rvert+\lvert b\rvert$; every $\VV$ qubit has
$\lvert a\rvert$ edges of type $X$ and $\lvert b\rvert$ of type $Z$, and every
$\CC$ qubit has $\lvert b\rvert$ of type $X$ and $\lvert a\rvert$ of type $Z$.
The graph is $\Delta$-regular, and $\Delta$ is the lower bound on the depth.

\section{Results}\label{sec:results}
\subsection{A counting identity}
For a schedule $\tau$ and a data qubit $q$, let
\[
I_q=\#\bigl\{(e,f): e\text{ an }X\text{-edge at }q,\ f\text{ a }Z\text{-edge at }q,\ \tau(e)<\tau(f)\bigr\}
\]
count the $X$-before-$Z$ orderings at $q$. Every such pair of edges joins $q$
to an $X$ ancilla $x$ and a $Z$ ancilla $z$ with $q\in N(x)\cap N(z)$, so
\begin{equation}\label{eq:doublecount}
\sum_{(x,z)}\ \sum_{q\in N(x)\cap N(z)}\one\bigl[\tau(x,q)<\tau(z,q)\bigr]
=\sum_q I_q .
\end{equation}
If $\tau$ is proper, every inner sum on the left is even, hence so is
$\sum_qI_q$.

\begin{lemma}\label{lem:count}
Let $\Delta\ge1$ and let $S\subseteq\{1,\dots,\Delta\}$. The number of pairs
$(s,t)$ with $s\in S$, $t\notin S$ and $s<t$ equals
\[
I(S)=\sum_{s\in S}(\Delta-s)-\binom{\lvert S\rvert}{2}.
\]
\end{lemma}
\begin{proof}
For $s\in S$ there are $\Delta-s$ layers above $s$, of which those in $S$ must
be removed. Summed over $s$, the removed pairs are the $\binom{\lvert S\rvert}{2}$
pairs $s<s'$ of elements of $S$.
\end{proof}

\begin{proposition}\label{prop:general}
Let $T$ be the Tanner graph of a CSS code with maximal degree $\Delta$, and
suppose that every $X$ ancilla and every data qubit has degree $\Delta$. Let
$m_X$ be the number of $X$ ancillas and $x_q$ the number of $X$-edges at the
data qubit $q$. Every schedule of depth $\Delta$ satisfies
\[
\sum_qI_q=m_X\binom{\Delta}{2}-\sum_q\binom{x_q}{2}.
\]
Consequently, if this number is odd, no proper schedule of depth $\Delta$
exists.
\end{proposition}
\begin{proof}
At depth $\Delta$, a vertex of degree $\Delta$ uses every layer exactly once.
At a data qubit $q$ the $X$-edges occupy a set $S_q$ of $x_q$ layers and the
$Z$-edges occupy its complement, so $I_q=I(S_q)$ by Lemma~\ref{lem:count}.
Summing over $q$, the first term of $I(S_q)$ becomes a sum over all
$X$-edges $e$ of $\Delta-\tau(e)$. Grouping the $X$-edges by their ancilla,
each of which uses every layer once, gives
$m_X\sum_{t=1}^{\Delta}(\Delta-t)=m_X\binom{\Delta}{2}$. The last statement
follows from \eqref{eq:doublecount}.
\end{proof}

\subsection{The main theorem}
\begin{theorem}\label{thm:main}
Let $G$ be a finite group of odd order $g$, and let $a,b\in\F[G]$ have odd
weights $\lvert a\rvert$ and $\lvert b\rvert$. The lifted product of the
$1\times1$ protographs $(a)$ and $(b)$ admits no proper schedule of depth
$\Delta=\lvert a\rvert+\lvert b\rvert$. Its minimal proper depth is $\Delta+1$,
attained by the three-band schedule of \cite{NRB2026}.
\end{theorem}
\begin{proof}
The Tanner graph \eqref{eq:edges} satisfies the hypotheses of
Proposition~\ref{prop:general} with $m_X=g$, with $g$ qubits $\VV_u$ having
$x_q=\lvert a\rvert$ and $g$ qubits $\CC_u$ having $x_q=\lvert b\rvert$. Hence
every schedule of depth $\Delta$ has
\[
\sum_qI_q=g\Bigl[\binom{\lvert a\rvert+\lvert b\rvert}{2}-\binom{\lvert a\rvert}{2}-\binom{\lvert b\rvert}{2}\Bigr]
=g\,\lvert a\rvert\,\lvert b\rvert ,
\]
which is odd. By Proposition~\ref{prop:general} no proper schedule of depth
$\Delta$ exists. The upper bound $\Delta_A+2\lceil\Delta_B/2\rceil=\Delta+1$ is
the depth of the three-band schedule of \cite{NRB2026}.
\end{proof}

The argument uses only the total parity, so it allows the cancellations
between different factorisations $u'=huk$ that the stricter per-summand
condition~(11) of \cite{NRB2026} forbids; it applies to every schedule of the
lifted graph, not only to schedules obtained by lifting.

For the three-band schedule we use the following explicit form, which the
checker verifies. Split $\supp b$ into a set $B_1$ of $\lceil\lvert b\rvert/2\rceil$
monomials and a set $B_2$ of the remaining ones. Layers
$1,\dots,\lceil\lvert b\rvert/2\rceil$ form the early band, the next
$\lvert a\rvert$ layers the middle band, and the last
$\lceil\lvert b\rvert/2\rceil$ layers the late band. The edges of $X_u$ and $Z_u$
coming from the same monomial $h\in\supp a$ share one middle layer. For
$k\in B_1$ the edge $X_u-\CC_{uk}$ is early and $Z_u-\VV_{uk^{-1}}$ late, and for
$k\in B_2$ the reverse. On every shared pair of qubits $\VV_{hu}$ and
$\CC_{uk}$ one of the two ancillas is first at both, so each pair contributes
an even number to \eqref{eq:proper}.

\subsection{Schedules uniform under the group}
The schedules of \cite{NRB2026} are obtained by lifting: the layer of an edge
depends only on its family in \eqref{eq:edges} and on its monomial $h$ or $k$,
not on $u$. Within this class the odd order is not needed.

\begin{proposition}\label{prop:uniform}
Let $G$ be any finite group and $a,b\in\F[G]$ of odd weights. No proper
schedule of depth $\Delta$ is uniform in the above sense.
\end{proposition}
\begin{proof}
Let $\mathcal A$, $\mathcal B$, $\mathcal B^*$ and $\mathcal A^*$ be the sets
of layers used by the four families $X-\VV$, $X-\CC$, $Z-\VV$ and $Z-\CC$. At
depth $\Delta$ each vertex uses every layer once. An $X$ ancilla shows that
$\mathcal A$ and $\mathcal B$ partition $\{1,\dots,\Delta\}$, a $\VV$ qubit that
$\mathcal A$ and $\mathcal B^*$ do, and a $\CC$ qubit that $\mathcal B$ and
$\mathcal A^*$ do. Hence $\mathcal B^*=\mathcal B$ and $\mathcal A^*=\mathcal A$.
Fix an $X$ ancilla $x$ and sum \eqref{eq:proper} over all $Z$ ancillas $z$. A
neighbour $\VV$ of $x$ reached in layer $s\in\mathcal A$ contributes the
number of its $Z$-edges in later layers, namely
$\#\{t\in\mathcal B: t>s\}$; a neighbour $\CC$ reached in layer $t\in\mathcal B$
contributes $\#\{s\in\mathcal A: s>t\}$. The total is the number of pairs
$(s,t)\in\mathcal A\times\mathcal B$, that is $\lvert a\rvert\,\lvert b\rvert$, which
is odd. Some pair $(x,z)$ therefore violates \eqref{eq:proper}.
\end{proof}

Summing the per-summand condition~(11) of \cite{NRB2026} over the
$\lvert a\rvert\,\lvert b\rvert$ pairs of monomials gives the same count, so the
obstruction also holds, a fortiori, in the more restricted class in which
that condition is imposed.

\subsection{Weight-two checks}
For $\lvert a\rvert=\lvert b\rvert=1$ the obstruction holds for every schedule
over every group. Each $X$ ancilla $x$ has one $\VV$ neighbour $q_1$ and one
$\CC$ neighbour $q_2$, and $q_1$ has a single $Z$ neighbour $z$. Because $x$
and $z$ commute, their overlap is even, so $z$ is also adjacent to $q_2$; each
connected component is a four-cycle $x-q_1-z-q_2-x$. A schedule of depth two
alternates its layers around the cycle: if $\tau(x,q_1)=c$ then
$\tau(z,q_1)=c'$, $\tau(z,q_2)=c$ and $\tau(x,q_2)=c'$, with $\{c,c'\}=\{1,2\}$.
The $X$ ancilla is first at exactly one of $q_1,q_2$, and \eqref{eq:proper}
fails.

\subsection{An application}
The bivariate bicycle code $[[90,8,10]]$ of \cite{Bravyi2024} has $\ell=15$,
$m=3$, $A=x^9+y+y^2$ and $B=1+x^2+x^7$, that is, $H_X=[A\;\;B]$ and
$H_Z=[B^{\mathsf T}\;\;A^{\mathsf T}]$ over $G=\mathbb{Z}_{15}\times\mathbb{Z}_3$.
It is the lifted product of two $1\times1$ protographs of weights $3$ and $3$
over a group of odd order $45$, with $\Delta=6$ and parity count
$45\cdot3\cdot3=405$. By Theorem~\ref{thm:main}, six layers of controlled-NOT
gates cannot measure all its checks in this model, and the seven layers of
the circuit of \cite{Bravyi2024} are optimal. The other codes of the table in
\cite{Bravyi2024} are built over groups of even order and are not covered.

\section{Even-order groups}\label{sec:even}
For a group of even order the total count $g\,\lvert a\rvert\,\lvert b\rvert$ is
even, and Proposition~\ref{prop:general} gives no obstruction. Exhaustive
search shows that the answer then depends on the code. For
$a=b=1+x+x^2$ over the cyclic group $\mathbb{Z}_6$, a CSS code with parameters
$[[12,4]]$, there is a proper schedule of depth $\Delta=6$, recorded as an
exact certificate with this entry and checked independently of the search
that produced it. The odd order in Theorem~\ref{thm:main} therefore cannot be
dropped, and by Proposition~\ref{prop:uniform} this schedule is not uniform
under the group. In contrast, for $a=1$ and $b=1+x+x^2$ over $\mathbb{Z}_4$ and
$\mathbb{Z}_6$ the search finds no proper schedule of depth $\Delta=4$. Whether
the Gross code, over a group of order $72$, admits depth $\Delta$ is not
decided here.

\section{Verification}\label{sec:verification}
The results of Section~\ref{sec:results} are proved analytically. The public
\href{https://github.com/jpbruneton/open_quantum_problems/blob/main/research/entries/syndrome-schedule-parity/checks/check_syndrome_parity.py}{Python checker},
the
\href{https://github.com/jpbruneton/open_quantum_problems/blob/main/research/entries/syndrome-schedule-parity/data/even_order_schedule.json}{depth-6 certificate over $\mathbb{Z}_6$}
and the
\href{https://github.com/jpbruneton/open_quantum_problems/blob/main/research/entries/syndrome-schedule-parity/data/verification.json}{recorded output}
form part of this entry. Over cyclic groups $\mathbb{Z}_g$ the checker builds
the graph \eqref{eq:edges} and tests the following. Lemma~\ref{lem:count} is
checked on all subsets for $\Delta\le9$. The identity
$\sum_qI_q=g\,\lvert a\rvert\,\lvert b\rvert$ is checked as an integer equality
on sixty random depth-$\Delta$ edge colourings of six codes, for which
properness is not imposed. Proposition~\ref{prop:uniform} is checked by
enumerating all uniform schedules of depth $\Delta$ on sixteen codes with odd
weights, over groups of odd and even order, with four controls of even weight
that do reach depth $\Delta$. An exhaustive backtracking search over all
schedules finds no proper schedule of depth $\Delta$ for $\lvert a\rvert=\lvert b\rvert=1$
over $\mathbb{Z}_g$ with $g\le8$, for the weights $(1,3)$ over
$\mathbb{Z}_g$ with $3\le g\le7$, and for the weights $(3,1)$ over
$\mathbb{Z}_3$ and $\mathbb{Z}_5$;
on the same codes the explicit three-band schedule of depth $\Delta+1$ is
proper. With the option \texttt{--extended} the search also exhausts
$a=b=1+x+x^2$ over $\mathbb{Z}_3$, about $3\times10^{8}$ search nodes and ten
minutes; the default run takes about ten seconds. Finally the checker
verifies the certificate and the parameters of the $[[12,4]]$ code, and the
parameters $[[90,8]]$, the degrees and the odd count $405$ for the bivariate
bicycle code of the previous section.

\section{Scope and limitations}\label{sec:scope}
Theorem~\ref{thm:main}, Propositions~\ref{prop:general} and~\ref{prop:uniform}
and the weight-two argument are proved analytically. The exhaustive searches
are independent computational confirmations on small cyclic groups, and the
depth-$\Delta$ schedule over $\mathbb{Z}_6$ is a computational result with an
exact, independently checked certificate. The matching upper bound
$\Delta+1$ is the construction of \cite{NRB2026}, which the checker tests on
small instances.

The results concern $1\times1$ protographs, circuits of controlled-NOT gates
with one ancilla per check, and the gate depth only. Larger protographs,
whose Tanner graphs need not be regular, flag qubits, ancilla reuse and other
gate sets are not treated, although Proposition~\ref{prop:general} applies to
any CSS code in which every $X$ ancilla and every data qubit has maximal
degree. For groups of even order with weights above one the minimal depth is
code-dependent, as Section~\ref{sec:even} shows, and remains open in general,
in particular for the Gross code.

On bibliographic priority, the obstruction was not found in
\cite{NRB2026}, which leaves the optimality of $\Delta+1$ as a numerical
suggestion, nor in the other works cited here. This is a statement about the
documented search only, not a claim of priority.

\clearpage
\begingroup\small
\begin{thebibliography}{99}
\setlength{\itemsep}{2pt}
\bibitem{NRB2026}
M. T. P. Nguyen, M. Rimbach-Russ and S. Bosco,
\emph{Disassembling qLDPC codes for depth-optimal parity-check circuits},
preprint (2026), version 2;
\href{https://arxiv.org/abs/2608.19917}{arXiv:2608.19917}.

\bibitem{TillichZemor2014}
J.-P. Tillich and G. Z\'emor,
\emph{Quantum LDPC codes with positive rate and minimum distance proportional
to the square root of the blocklength},
IEEE Trans. Inf. Theory \textbf{60}, 1193 (2014);
\href{https://doi.org/10.1109/TIT.2013.2292061}{doi:10.1109/TIT.2013.2292061};
\href{https://arxiv.org/abs/0903.0566}{arXiv:0903.0566}.

\bibitem{PanteleevKalachev2021}
P. Panteleev and G. Kalachev,
\emph{Degenerate quantum LDPC codes with good finite length performance},
Quantum \textbf{5}, 585 (2021);
\href{https://doi.org/10.22331/q-2021-11-22-585}{doi:10.22331/q-2021-11-22-585};
\href{https://arxiv.org/abs/1904.02703}{arXiv:1904.02703}.

\bibitem{LinPryadko2023}
H.-K. Lin and L. P. Pryadko,
\emph{Quantum two-block group algebra codes},
preprint (2023);
\href{https://arxiv.org/abs/2306.16400}{arXiv:2306.16400}.

\bibitem{Bravyi2024}
S. Bravyi, A. W. Cross, J. M. Gambetta, D. Maslov, P. Rall and T. J. Yoder,
\emph{High-threshold and low-overhead fault-tolerant quantum memory},
Nature \textbf{627}, 778 (2024);
\href{https://doi.org/10.1038/s41586-024-07107-7}{doi:10.1038/s41586-024-07107-7};
\href{https://arxiv.org/abs/2308.07915}{arXiv:2308.07915}.

\bibitem{TDB2022}
M. A. Tremblay, N. Delfosse and M. E. Beverland,
\emph{Constant-overhead quantum error correction with thin planar
connectivity},
Phys. Rev. Lett. \textbf{129}, 050504 (2022);
\href{https://doi.org/10.1103/PhysRevLett.129.050504}{doi:10.1103/PhysRevLett.129.050504};
\href{https://arxiv.org/abs/2109.14609}{arXiv:2109.14609}.
\end{thebibliography}
\endgroup
\end{document}
